% ======================================================
%           《 Ac's Life OS：人生操作系统维护协议 》
% ======================================================
% 文件属性：现行宪法 / 外部大脑 / 进化日志
% 当前版本：v2026.A
% 最后修订：\today
% ======================================================
%
% 【 核心维护原则：正文即现行宪法 】 
% 1. 实时性：正文永远只保留你此时此刻“正在执行”的最佳实践。
% 2. 纯粹性：一旦某项工具或习惯被弃用，立即移至 [Archive.tex]，保持正文精简。
% 3. 场景化：记录工具时，必须包含“触发场景”（如：当焦虑时、当需要深度构思时）。
%
% 【 迭代更新指南：记什么 & 不记什么 】
% - ✅ 记【认知升级】：为什么要从工具A换到工具B？当时的痛点是什么？
% - ✅ 记【判定逻辑】：比如“B站视频3次未看完即删除”的心理阈值。
% - ❌ 不记【琐碎修改】：修正错别字、调整颜色、更改字体等无需记录。
% - ❌ 不记【过去时】：不要在正文中使用“我以前习惯...”这样的表达。
%
% 【 第二大脑 (BASB) 实践记录 】
% - C (Capture)：抓取有益论述（YouTube/B站/书籍）。
% - O (Organize)：模块化管理（Philosophy, Toolkit, Maintenance）。
% - D (Distill)：提炼核心 SOP（清理准则、仪式感流程）。
% - E (Express)：将输入转化为这一份可执行的工作录。
%
% 【 数字化代谢协议 (Metabolic SOP) 】
% - 每日：睡前清理手机截图，未入 Notion 库者一律销毁。
% - 每周：周日晚 21:00 执行收藏夹“生死判决”，超过 7 天未处理则删除。
% - 每季：审视 [Archive.tex]，观察自己的思维化石，复盘进化路径。
%
% ======================================================
%           仪式感触发：清空桌面，启动深度思考模式
% ======================================================
%npm run push-log   是可以自动到github    ，ctrl加~按键启动终端
%调用  让localhost能够加载出来命令：npm run dev（虽然这个应该放到最终的readme之中）
\documentclass[11pt, a4paper, oneside]{book} 

% 调用我们自定义的宏包
\usepackage{workingonon}

% ======================================================
% 5. 文档开始
% ======================================================
\begin{document}

% --- 封面 ---
\begin{titlepage}  
    \centering
    \vspace*{5cm}
    {\Huge \bfseries Ac 的工作录与生活协议} \\[1cm]
    {\large \itshape --- “为了更清醒地工作，更诗意地生活” ---} \\[4cm]
    \textbf{持有人：} Ac (阿曹··奥特兄弟) \\
    \textbf{最后更新：} \today \\
    \vspace{2cm}
    % 这里可以放置你的 Logo 占位
    \includegraphics[width=13cm]{pp1.png}
\end{titlepage}

\tableofcontents
\newpage

% ======================================================
% 第一部分：输入与沉淀 (Media & Insights)
% ======================================================
\part*{零：总则}
\begin{dododo}{v2.1 - 2026.03.03}
    

1.主线贯穿支线，固定SOP不赘述，分清场合，逼自己花时间。

2.不要吝啬出去走走的时间

3.不必过于纠结工具的选择，重要的是使用它们的场景和目的。

4.定期回顾和调整工具使用习惯，保持灵活性和适应性。要流动起来，而非一了百了，这才是复盘和整理的意义

比如我自己的工具也是从微信聊天框变成滴答清单然后到苹果提醒事项与日历然后一点一点分出来的，他其实是一个螺旋渐进的过程，工具的迭代也是一个不断试错和优化的过程。

5.保持对新工具的开放态度（尤其是ai），但也要有选择性地尝试。

6.助人者自助，对自己有用。

7.有空记得雕花，渲染页面。

8.最小执行单位。

\end{dododo}

\part{输入与沉淀}

\chapter{影音与社交媒体论述}
\section*{1.0  随想}

\begin{flowing}{v2.1 - 2026.2.27夜}
    我现在觉得数码产品越做越有意思。比如苹果的快捷指令（给出了很多能编程的小组件，本质是很多小接口）；再比如notion，他们的核心设计思路都是把功能做到最简，包括vscode。留了一个可以编程的功能给用户，让用户自己发挥。相当于是视觉编程，善于学习和探索的聪明人可以把苹果手机甚至别的所有手机玩出花来，不然就只是纯粹的显示器和打字机。也有点类似于scratch小黄猫。
这其实有点考验人对手机的理解。至少苹果公司在引导人去利用工具，而不是被工具异化，尽管有一定门槛。
\end{flowing}

\section{Bilibili：从消遣到课堂}
这里记录你对 B 站的使用逻辑。例如：只看技术类与人文类，拒绝情绪化短视频。

\begin{updatelog}{v2.1 - 2026.02.26}
    删除了默认收藏夹里的10000多个视频，设定了“3次未看完即删除”的清理规则。
    顺带删除了Ac微信对话框里的信息。重新构建收藏夹认知：少，即是多。
\end{updatelog}

\section{Youtube 与 NotebookLM 的联动}
论述如何将 Youtube 的长视频内容喂给 NotebookLM 进行二次消化。

\section{CT推的ima知识库，但我还没试过}
RAG或者在线知识库？？
%====================================
% 第二部分：工具阵列 (Tools Arsenal)
% ======================================================
\part{工具阵列}

\chapter{模拟工具：仪式感与手写}
\section{纸质笔记本与定制信纸}

描述你在什么场景下使用那份带有 Logo 的仪式化信纸。是用于给重要的思绪定调，还是用于某种特定的复盘？（虚线提示本句）

四象限倾向于用纸和笔来描述

\section{iPad 手写板：动态草稿}

\begin{updatelog}{v2.1 - 2026.02}
   
    花了一点小钱买了notability，goodnotes，以及无边记的思维导图功能开始解锁，放弃了一年68的万兴pdf，但要注意扫描全能王工具的平替。
    要把gn和nb的功能解锁好
    
\end{updatelog}

\chapter{数字工具：结构化大脑}

\section{{$\mathfrak{Latex}$}与多种功能}
\begin{flowing}{v2.1 - 2026.03.18}
    有想法希望用latex可以做课程笔记（down下来了一个非常好的模板），可以有$input$或者$include$插入命令，python制作图与表。甚至教材也可以编写。
\end{flowing}

\begin{updatelog}{v2.11 - 2026.03.19}
    工作std：把宏包等导言切出去，正文专心编辑与排版
\end{updatelog}

\begin{dododo}{v2.1 - 2026.03.19}
    功能一定要捋干净，要分模块管理思想和生活！！
   
\end{dododo}

\section{苹果内置备忘录：快速捕捉}
\begin{updatelog}{v2.1 - 2026.02}
    结合快捷指令，他现在可以支持剪贴板或者语音或者文字的灵感快速录入
\end{updatelog}

\section{Notion：万物索引}
\begin{flowing}{v0 - 2026.03.03}
Notion 作为一个“编程式”的工具，适合用来构建复杂的知识体系和项目管理。它的数据库功能可以帮助你建立起一个动态更新的“第二大脑”，将各种信息进行模块化管理。

\large\textbf{差别：}

Notion的思维模式是延时满足，要把很多信息在脑子里过很多遍，然后写到notion里，相当于临门一脚，相信大脑加工的力量。

而微信的思维模式是想一出是一出，可以记录完整，但是清理和复盘的成本较高。尽管想一出是一出比较符合人探索的心情，而且不适合给到灵感速记，这个规划不能算灵感。
\end{flowing}
\begin{updatelog}{v2.1 - 2026.02.28}
    
    将原来放于提醒事项中的部分内容放给notion，使得归档与提醒分得开一些。利用其可以“编程”的属性，但还没有渲染界面，买了一些教程，可以学一学。
    intention：以后学习规划就在notion里单开一页
\end{updatelog}

\begin{updatelog}{v2.11 - 2026.03.03}
    
    建立学习规划页面，目前主要是当日小任务，主要是复习预习与应付作业，后续会延伸出备考环节。
    
\end{updatelog}

\begin{updatelog}{v2.12 - 2026.03.04}
    
   周复盘，建立一周支线sop，四个象限内容建立成功。打卡页面（48*3）的点击加一功能做出）
    
\end{updatelog}

\begin{updatelog}{v2.13 - 2026.03.20}
    
   Mermaid+mindmap指令，可以制作很好的思维导图。删掉了原来的思维导图制作软件
   刚看了一下他事实上有更多的编程模块，比如python，latex……，简直了大大的轻量编辑器。
    
\end{updatelog}

\begin{updatelog}{v2.1？ - }
     
   联不联claude呢。

\end{updatelog}

\section{Calflow}

\begin{updatelog}{v2.1 - 2026.02.27}
    与ios联动的时间管理与统计工具，可视化功能好用，可以进一步优化时间分配与效率分析。
\end{updatelog}

\section{收藏夹定期清理协议 (SOP)}
\begin{itemize}
    \item \textbf{执行时间：} 每周日晚 21:00。
    \item \textbf{清理对象：} B 站“稍后再看”、手机图库截图、浏览器书签。
    \item \textbf{判定标准：} 超过 7 天未处理且无长效价值的内容，直接删除。
\end{itemize}

\chapter{AI 助手：未来的工作伙伴}

\section{deepseek:高情商学术不端垃圾}

\begin{updatelog}{v2.1 - 2026.03.03}
     
    不更新之前基本上不会用，打算删了

\end{updatelog}

\begin{updatelog}{v2.11 - 2026.03.07}
     
    api接入claude，充了20r。

\end{updatelog}

\begin{updatelog}{v2.12 - 2026.03.11}
     
    api接出，贝利亚黄昏与泽塔奥特曼，牛逼的模拟聊天工具。

\end{updatelog}

\section{Gemini 3.0 pro}
注意gemini的高阶用法，插件，权益
\begin{updatelog}{v2.10 - 2026.02.14}

    买号充会员，他给我了很多很好的想法，包括本文档也是Gemini给我的框架。但是他对上下文的黏性太大了，即便开启新对话他还是记得我和他说的所有上下文。

\end{updatelog}

\begin{updatelog}{v2.11 - 2026.03.14}
     
   在调教物理实验报告模板上，对.0的处理，最后的结果是换用python生图，latex主做排版。

   这启示我们对于ai提问上要给出更准确的方向，比如是换工具还是深刻纠错等。

\end{updatelog}

\section{豆包手机版}
\begin{inspiration}{v2.10 - 2026.03.07}
     
  在教人干活的方面，豆包可以做到出答案快准且命中率高，废话少，简练。比Gemini更加不让人飙血压。

\end{inspiration}

\begin{updatelog}{v2.11 - 2026.03.11}
     
    api接出，变成了桌面宠物美弗拉斯星人（有问必答星人），但是回答速度较deepseek更慢，回答内容更丰富。

\end{updatelog}

\section{grok4.2}
据说能够多agent讨论

\section{GPT}
to be continued……
% ==================

\begin{updatelog}{v2.11 - 2026.03.03}
    
    准备和zny一起用，买会员，研究一下家庭组（待做）

\end{updatelog}

\section{Claude code}
\begin{updatelog}{v2.1 - 2026.03.06}
    
    配置好环境，vscode或者powershell都可以用，接入了豆包和ds的模型，可以直接做网页，api接入后的可操作性比较强，算是工作的利器。
    缺点就是api要钱，虽然不贵。有待后续各种工作与项目的实际开发，以及与notion等的联动。
    记得学学agent skill之类的

\end{updatelog}

% ======================================================
% 第三部分：年度回顾 (Optional)
% ======================================================
\part{年度演进}
\begin{sum1}{v2.1 - 2026.03.19}
\Large\textbf{关键词（26.3.3）：迭代与优化}
\end{sum1}

\begin{sum2}{v2.1 - 2026.03.19}
\Large\textbf{关键词（26.3.3）：轻量化与可视化}

\end{sum2}

\begin{sum3}{v2.1 - 2026.03.19}
\Large\textbf{关键词（26.3.3）：结构化与模块化}\\
（notion，latex，日程，calflow etc）
\end{sum3}

\end{document}